\chapter[22]{Tools Used in Proof of Pseudolocality}
\label{ch:chow-iii-tools-pseudolocality}
\dcref{ch22}
\source{chow-et-al-ricci-flow-part-iii:page-0204}

Set \(\mathcal M_\alpha=\{(x,t):|\Rm|(x,t)>\alpha/t\}\).

\section{Point picking}

\begin{definition}[\(\alpha\)-small and \(\alpha\)-large curvature points]
\label{def:chow-iii-alpha-small-large-curvature}
\source{chow-et-al-ricci-flow-part-iii:page-0205}
\dcref{ch22:22.1}
\group{Point picking}

For \(\alpha>0\), a space-time point \((x,t)\), \(t>0\), is
\emph{\(\alpha\)-small} if \(|\Rm|(x,t)\le\alpha/t\), and is
\emph{\(\alpha\)-large} if \(|\Rm|(x,t)>\alpha/t\).  The condition is
invariant under the parabolic scaling \(g_C(\tau)=Cg(\tau/C)\).
\end{definition}

\begin{lemma}[Large curvature point with local curvature control]
\label{lem:chow-iii-large-curvature-local-control}
\source{chow-et-al-ricci-flow-part-iii:page-0205,page-0206}
\dcref{ch22:22.3}
\group{Point picking}

Let \(\alpha>0\), \(\epsilon\in(0,1)\), and let a bounded-curvature flow on
\([0,T]\) contain \((x_0,t_0)\in\mathcal M\times[T/2,T]\) with
\[
 |\Rm|(x_0,t_0)>\max\{\alpha/t_0,8/(T\epsilon)\}.
\]
Then there is an \(\alpha\)-large point \((\hat x_0,\hat t_0)\),
\(\hat t_0\in[T/4,T]\), with \(\hat Q=|\Rm|(\hat x_0,\hat t_0)\) at least
\(8/(T\epsilon)\), such that
\[
 |\Rm|(x,t)\le2\hat Q
\]
on \(P(\hat x_0,\hat t_0,(epsilon\hat Q)^{-1/2},-(\epsilon\hat Q)^{-1})\).
\end{lemma}
\begin{proof}
\sketch If the initial point does not have the asserted property, choose a point in
its backward parabolic cylinder with curvature greater than twice the current
value and iterate.  The curvature grows geometrically while each time step is
at most half the preceding reciprocal curvature scale, so the times remain in
\([T/4,T]\).  Bounded curvature makes the iteration terminate, proving the
claim.
\end{proof}

\begin{claim}[Picking a well-chosen large curvature point]
\label{clm:chow-iii-point-picking-claim-one}
\source{chow-et-al-ricci-flow-part-iii:page-0207,page-0208}
\uses{def:chow-iii-alpha-small-large-curvature}
\dcref{ch22:Claim1}
\group{Point picking}


Let \(\epsilon_0,\alpha,A>0\), \(\epsilon\le\epsilon_0\), and suppose that
\((x_1,t_1)\) satisfies
\[
 x_1\in B_{g(t_1)}(x_0,\epsilon_0),\qquad
 |\Rm|(x_1,t_1)>\alpha/t_1+\epsilon_0^{-2}.
\]
There is an \(\alpha\)-large point \((\bar x,\bar t)\), with
\(\bar t\in(0,\epsilon^2]\) and
\(\bar x\in B_{g(\bar t)}(x_0,(2A+1)\epsilon_0)\), such that every
\((x,t)\in\mathcal M_\alpha\) with \(t\le\bar t\) and
\[
 d_{g(t)}(x,x_0)\le d_{g(\bar t)}(\bar x,x_0)+A|\Rm|(\bar x,\bar t)^{-1/2}
\]
satisfies \(|\Rm|(x,t)\le4|\Rm|(\bar x,\bar t)\).
\end{claim}
\begin{proof}
\sketch Starting at \((x_1,t_1)\), choose a later point only when the asserted bound
fails, increasing curvature by a factor greater than four and increasing the
distance by at most the corresponding curvature scale.  The geometric series
of scales keeps all points in the stated ball.  A nonterminating sequence would
contradict bounded curvature on the compact space-time tube, so the process
terminates.
\end{proof}

\begin{remark}[Compactness suffices in Claim 1]
\label{rem:chow-iii-point-picking-compactness}
\source{chow-et-al-ricci-flow-part-iii:page-0208}
\uses{clm:chow-iii-point-picking-claim-one}
\dcref{ch22:22.4}
\group{Point picking}


The proof of Claim 1 only needs compactness of the closed balls
\(\overline B_{g(t)}(x_0,(2A+1)\epsilon_0)\) for each time, rather than global
completeness.
\end{remark}

\begin{lemma}[Point picking on an arbitrary space-time set]
\label{lem:chow-iii-point-picking-subset}
\source{chow-et-al-ricci-flow-part-iii:page-0209}
\uses{clm:chow-iii-point-picking-claim-one}
\dcref{ch22:22.5}
\group{Point picking}


If \(\mathcal S\subset\mathcal M\times(0,\epsilon^2]\) contains a point within
\(\epsilon_0\) of \(x_0\) whose curvature exceeds \(\epsilon_0^{-2}\), then
there is \((\bar x,\bar t)\in\mathcal S\), with \(\bar t\le t_1\) and
\(\bar x\in B_{g(\bar t)}(x_0,(2A+1)\epsilon_0)\), whose curvature exceeds
\(\epsilon_0^{-2}\) and dominates by a factor four every point of \(\mathcal S\)
not farther from \(x_0\) by more than \(A|\Rm|(\bar x,\bar t)^{-1/2}\).
\end{lemma}

\begin{claim}[The well-chosen point has backward parabolic curvature control]
\label{clm:chow-iii-point-picking-claim-two}
\source{chow-et-al-ricci-flow-part-iii:page-0209,page-0210,page-0211,page-0212}
\uses{clm:chow-iii-point-picking-claim-one, lem:chow-iii-point-picking-local-cylinder}
\dcref{ch22:Claim2}
\group{Point picking}


Assume \(A\ge1\) and \(\alpha<1/(13(n-1)\sqrt n)\), and let
\(\bar Q=|\Rm|(\bar x,\bar t)\) be supplied by Claim 1.  Then
\[
 |\Rm|(x,t)\le4\bar Q
\]
whenever \(\bar t-\alpha/(2\bar Q)\le t\le\bar t\) and
\(d_{g(t)}(x,\bar x)\le A\bar Q^{-1/2}/10\).  The same conclusion holds
with radius \(A\bar Q^{-1/2}\) as noted in Remark 22.7.
\end{claim}
\begin{proof}
\sketch At points which are not \(\alpha\)-large the bound follows from
\(|\Rm|\le\alpha/t<2\bar Q\).  For \(\alpha\)-large points, the preceding
lemma places the relevant parabolic cylinder inside the distance-controlled
region of Claim 1, giving the factor-four estimate.
\end{proof}

\begin{lemma}[Cylinder inclusion for Claim 2]
\label{lem:chow-iii-point-picking-local-cylinder}
\source{chow-et-al-ricci-flow-part-iii:page-0210,page-0211,page-0212}
\uses{clm:chow-iii-point-picking-claim-one}
\dcref{ch22:22.8}
\group{Point picking}


If \((x,t)\in\mathcal M_\alpha\) satisfies
\(\bar t-\alpha/(2\bar Q)\le t\le\bar t\) and
\(d_{g(t)}(x,\bar x)\le A\bar Q^{-1/2}/10\), then
\[
 d_{g(t)}(x,x_0)\le d_{g(t)}(\bar x,x_0)+A\bar Q^{-1/2}.
\]
Consequently the entire cylinder in Claim 2 lies in the region controlled by
Claim 1.
\end{lemma}
\begin{proof}
\sketch The triangle inequality gives the inclusion at time \(\bar t\).  If it first
failed at an earlier time, apply the changing-distances estimate on the
intervening interval, using the curvature bound from Claim 1 and the
non-large-point estimate.  The resulting distance decrease is strictly less
than the available \(9A\bar Q^{-1/2}/10\), a contradiction.
\end{proof}

\section{Heat kernels under Cheeger--Gromov limits}

\begin{lemma}[Convergence of heat kernels under Cheeger--Gromov limits]
\label{lem:chow-iii-heat-kernel-cheeger-gromov}
\source{chow-et-al-ricci-flow-part-iii:page-0212,page-0213,page-0214,page-0215,page-0216,page-0217}
\dcref{ch22:22.9}
\group{Heat kernels}

Let pointed complete backward Ricci flows \((\mathcal M_i,g_i(\tau),x_i)\),
\(\tau\in[0,\omega]\), with bounded curvature converge smoothly in the pointed
Cheeger--Gromov sense to \((\mathcal M_\infty,g_\infty(\tau),x_\infty)\).  If
\(H_i=(4\pi\tau)^{-n/2}e^{-f_i}\) are the adjoint heat kernels based at
\((x_i,0)\), and \(v_i\) their Perelman Harnack quantities, then after a
subsequence the pullbacks converge smoothly on compact subsets to
\(H_\infty=(4\pi\tau)^{-n/2}e^{-f_\infty}>0\), \(f_\infty\), and \(v_\infty\le0\).
Moreover
\[
 (\partial_\tau-\Delta_{g_\infty}+R_{g_\infty})H_\infty=0,
\]
and
\[
 (\partial_\tau-\Delta_{g_\infty}+R_{g_\infty})v_\infty
 =-2\tau\left|\Rc+\nabla^2f_\infty-\frac{g_\infty}{2\tau}\right|^2H_\infty.
\]
\end{lemma}
\begin{proof}
\sketch Local volume comparison and conservation of heat-kernel mass give uniform
integral bounds.  A Li--Yau estimate (or the parabolic mean-value inequality)
then gives local upper bounds away from \(\tau=0\); Dirichlet comparison with
a small constant-curvature ball gives a uniform positive lower bound at the
base points.  Interior derivative estimates and a diagonal subsequence yield
smooth convergence.  Passing to the limit in the equations and in
Perelman's inequality proves the stated identities.
\end{proof}

\begin{remark}[The limit kernel is expected to be minimal]
\label{rem:chow-iii-heat-kernel-minimality}
\source{chow-et-al-ricci-flow-part-iii:page-0218}
\uses{lem:chow-iii-heat-kernel-cheeger-gromov}
\dcref{ch22:22.11}
\group{Heat kernels}


The identification of \(H_\infty\) with the minimal positive fundamental
solution requires convergence to \(\delta_{x_\infty}\) as
\(\tau\downarrow0\).
\end{remark}

\section{Local entropy}

\begin{lemma}[Time-dependent cutoff subsolution]
\label{lem:chow-iii-time-dependent-cutoff}
\source{chow-et-al-ricci-flow-part-iii:page-0218,page-0219,page-0220}
\dcref{ch22:22.12}
\group{Local entropy}

Let \(\phi\colon\mathbb R\to[0,1]\) be smooth, equal to one on
\((-infty,1]\), zero on \([2,\infty)\), and satisfy
\((\phi')^2\le10\phi\), \(\phi''\ge-10\phi\).  For
\[
 h(x,t)=\phi\!\left(\frac{d_{g(t)}(x,x_0)+200n\sqrt t}{b}\right),
\]
assume \(\alpha\le1\), \(|\Rm|\le\alpha/t+2/\epsilon_0^2\) on
\(B_{g(t)}(x_0,\epsilon_0)\), and \(b\ge5\epsilon_0A\), \(A\ge41n\).  Then,
in the weak sense,
\[
 (\partial_t-\Delta_{g(t)})h\le\frac{10}{b^2}h.
\]
\end{lemma}
\begin{proof}
\sketch At points where \(\phi'\ne0\), the distance is at least \(b-200n\sqrt t\),
so the distance-Laplacian comparison theorem applies with
\(K=\alpha/t+2/\epsilon_0^2\).  It gives
\((\partial_t-\Delta)d_{g(t)}(x,x_0)\ge-100n/\sqrt t\), which cancels the
positive derivative of \(200n\sqrt t\).  The bound on \(-\phi''\) and
\(|\nabla d|=1\) yield the result.
\end{proof}

\begin{lemma}[Upper bound for local entropy at time zero]
\label{lem:chow-iii-local-entropy-upper-bound}
\source{chow-et-al-ricci-flow-part-iii:page-0220,page-0221,page-0222,page-0223,page-0224,page-0225,page-0226,page-0227}
\uses{lem:chow-iii-time-dependent-cutoff, clm:chow-iii-point-picking-claim-two}
\dcref{ch22:22.13}
\group{Local entropy}


Let \((\mathcal M^n,g(t),x_0)\), \(t\in[0,\epsilon^2]\), be a complete bounded-
curvature Ricci flow satisfying
\[
 |\Rm|(x,t)\le\alpha/t+2\epsilon_0^{-2}
\quad\text{when }d_{g(t)}(x,x_0)\le\epsilon_0,
\]
with \(\alpha\le1\).  Suppose \(A\ge67n\), \((\bar x,\bar t)\) is a point
with \(\bar Q=|\Rm|(\bar x,\bar t)>\alpha/\bar t\) and the distance and
curvature-control conclusions of Claim 1, and suppose that for some
\(\tilde t\in[\bar t-\alpha/(2\bar Q),\bar t]\),
\[
 \int_{B_{g(\tilde t)}(\bar x,\sqrt{\bar t-\tilde t})}v\,d\mu_{g(\tilde t)}\le-\beta_1.
\]
For the cutoff \(h\) above with \(b=10\epsilon_0A\),
\[
 \int_{\mathcal M}v h\,d\mu_{g(0)}
 \le-\beta_1(1-A^{-2}).
\]
If also \(A\le(20\epsilon_0)^{-1}\) and \(R(\cdot,0)\ge-1\) on
\(B_{g(0)}(x_0,1)\), then for \(\widetilde H=Hh\) and
\(\widetilde f=f-\log h\),
\[
 \int\bigl(-\tilde t|\nabla\widetilde f|^2-\widetilde f+n\bigr)\widetilde H\,d\mu_{g(0)}
 \ge \beta_1(1-A^{-2})-A^{-2}-\epsilon^2,
\]
and \(\int\widetilde H\,d\mu_{g(0)}\le1\).
\end{lemma}
\begin{proof}
\sketch The adjoint equation and \(\Box^*v\le0\), together with the cutoff
subsolution, give a Gronwall inequality for \(\int(-v)h\).  Integrating from
\(\tilde t\) to zero and using the distance estimate places the negative-entropy
ball inside the region where \(h=1\).  This proves the first inequality.  Writing
\(f=\widetilde f+\log h\), integration by parts reduces the second inequality
to bounds for \(-Rh\), \(|\nabla h|^2/h\), and \(-h\log h\).  The curvature lower bound,
\((\phi')^2\le10\phi\), and a second cutoff heat-kernel comparison bound these
terms by \(\epsilon^2\), \((10A^2)^{-1}\), and \(1-(9/10)A^{-2}\), respectively.
\end{proof}

\section{Logarithmic Sobolev inequalities}

\begin{definition}[Logarithmic Sobolev inequality]
\label{def:chow-iii-log-sobolev}
\source{chow-et-al-ricci-flow-part-iii:page-0227}
\dcref{ch22:22.14}
\group{Logarithmic Sobolev}

A Riemannian manifold \((\mathcal M^n,g)\) satisfies a logarithmic Sobolev
inequality if some \(C_P>-\infty\) satisfies
\[
 \int_{\mathcal M}\left(\frac12|\nabla f|^2+f-n\right)u\,d\mu\ge C_P
\]
for every \(u=(2\pi)^{-n/2}e^{-f}\) with \(\int u\,d\mu=1\).  Equivalently,
for \(\psi\in W^{1,2}\),
\[
 \int(2|\nabla\psi|^2-\psi^2\log\psi^2)
 +\log\!\left(\int\psi^2\right)\int\psi^2
 \ge\left(\tfrac n2\log(2\pi)+n+C_P\right)\int\psi^2.
\]
\end{definition}

\begin{theorem}[Euclidean logarithmic Sobolev inequality]
\label{thm:chow-iii-euclidean-log-sobolev}
\source{chow-et-al-ricci-flow-part-iii:page-0227}
\uses{def:chow-iii-log-sobolev}
\dcref{ch22:22.15}
\group{Logarithmic Sobolev}


For every \(\varphi\in W^{1,2}(\mathbb R^n)\),
\[
 \int_{\mathbb R^n}(2|\nabla\varphi|^2-\varphi^2\log\varphi^2)
 +\log\!\left(\int\varphi^2\right)\int\varphi^2
 \ge\left(\tfrac n2\log(2\pi)+n\right)\int\varphi^2.
\]
\end{theorem}

\begin{theorem}[Logarithmic Sobolev via isoperimetric inequality]
\label{thm:chow-iii-log-sobolev-isoperimetric}
\source{chow-et-al-ricci-flow-part-iii:page-0228,page-0229,page-0230,page-0231,page-0232}
\uses{thm:chow-iii-euclidean-log-sobolev}
\dcref{ch22:22.16}
\group{Logarithmic Sobolev}


If a complete \((\mathcal M^n,g)\) satisfies
\[
 \Area(\partial\Omega)^n\ge I_n\Vol(\Omega)^{n-1}
\]
for compact domains with \(C^1\) boundary, then every \(\psi\in W^{1,2}(\mathcal M)\)
satisfies
\[
 \int(2|\nabla\psi|^2-\psi^2\log\psi^2)
 +\log\!\left(\int\psi^2\right)\int\psi^2
 \ge\left(s_n+\log(I_n/c_n)\right)\int\psi^2,
\]
where \(s_n=\frac n2\log(2\pi)+n\) and \(c_n=n^n\omega_n\).
\end{theorem}
\begin{proof}
\sketch Rescale so that the isoperimetric constant is \(c_n\).  Approximate by a
nonnegative compactly supported \(C^1\) function and perform spherical
symmetrization: the radial rearrangement \(h\) has the same distribution
function, while the isoperimetric inequality and co-area formula imply
\[
 \int|\nabla h|^2\,d\mu_{\mathbb R^n}
 \le(c_n/I_n)^{2/n}\int|\nabla\psi|^2\,d\mu_g.
\]
The distribution identity preserves the logarithmic and \(L^2\) terms.  Apply
the Euclidean inequality and undo the scaling.
\end{proof}

\begin{remark}[Co-area interpretation]
\label{rem:chow-iii-log-sobolev-coarea}
\source{chow-et-al-ricci-flow-part-iii:page-0228,page-0229,page-0230}
\uses{thm:chow-iii-log-sobolev-isoperimetric}
\dcref{ch22:22.17}
\group{Logarithmic Sobolev}


The link between isoperimetric and Sobolev inequalities is furnished by the
co-area formula applied to the level sets of a function; spherical
symmetrization makes each comparison level set a Euclidean sphere.
\end{remark}

\begin{theorem}[Local logarithmic Sobolev inequality]
\label{thm:chow-iii-local-log-sobolev}
\source{chow-et-al-ricci-flow-part-iii:page-0232}
\uses{thm:chow-iii-log-sobolev-isoperimetric}
\dcref{ch22:22.19}
\group{Logarithmic Sobolev}


Let \(\overline B(x_0,\rho)\) be compact and suppose the isoperimetric bound
\(\Area(\partial\Omega)^n\ge I_n\Vol(\Omega)^{n-1}\) holds for compact domains
inside \(B(x_0,\rho)\).  Every compactly supported \(C^1\) function in this ball
satisfies
\[
 \int(2|\nabla\psi|^2-\psi^2\log\psi^2)
 +\log\!\left(\int\psi^2\right)\int\psi^2
 \ge\left(\epsilon_n+\log(I_n/c_n)\right)\int\psi^2,
\]
with the dimensional constant \(\epsilon_n\) from the source normalization.
\end{theorem}
\begin{proof}
\sketch Apply Theorem~\ref{thm:chow-iii-log-sobolev-isoperimetric} to the
positive superlevel sets of \(\psi\), all of which remain inside the ball.
\end{proof}

\begin{theorem}[Backward uniqueness for solutions of Ricci flow]
\label{thm:chow-iii-backward-uniqueness}
\source{chow-et-al-ricci-flow-part-iii:page-0232,page-0233}
\dcref{ch22:22.20}
\group{Uniqueness}

If \(g_1(t)\) and \(g_2(t)\) are complete bounded-curvature Ricci flows on
\([0,T]\) with \(g_1(T)=g_2(T)\), then \(g_1(t)=g_2(t)\) for every
\(t\in[0,T]\).
\end{theorem}

\begin{corollary}[The isometry group is preserved under Ricci flow]
\label{cor:chow-iii-isometry-group-preserved}
\source{chow-et-al-ricci-flow-part-iii:page-0233}
\uses{thm:chow-iii-backward-uniqueness}
\dcref{ch22:22.21}
\group{Uniqueness}


For a complete bounded-curvature Ricci flow,
\[
 \operatorname{Isom}(g(t))=\operatorname{Isom}(g(0))
\]
for every \(t>0\).
\end{corollary}
\begin{proof}
\sketch Backward uniqueness gives \(\operatorname{Isom}(g(t))\subseteq\operatorname{Isom}(g(0))\);
forward uniqueness gives the reverse inclusion.
\end{proof}

\begin{remark}[Criterion for canonical Ricci flow with surgery]
\label{rem:chow-iii-canonical-surgery-criterion}
\source{chow-et-al-ricci-flow-part-iii:page-0233,page-0234}
\dcref{ch22:22.22}
\group{Uniqueness}

The chapter proposes that a canonical Ricci flow with surgery on closed
three-manifolds should have the property that an isometry between open subsets
at one time determines an isometry of the universal covers at every time,
including across surgery times.
\end{remark}
